\section{Related work}

The closest identification and efficiency analysis is that of \citet[Lemma~1.1 under Assumptions~1--3, Theorem~2.1, and Sections~3--4]{KallusMao2025}. Their experiment observes baseline covariates, treatment, and surrogate outcomes for all units, together with primary outcomes for a labeled subset. Treatment unconfoundedness, missingness at random conditional on covariates, treatment, and surrogates, and the corresponding overlap conditions identify the population average treatment effect. Theorem~2.1 characterizes the semiparametric efficiency bound; Sections~3--4 develop estimation and asymptotic inference under nuisance convergence and moment conditions, including a regime with a dominant unlabeled sample. Surrogates supply predictive information about the primary outcome within this identification framework. The present analysis specializes treatment assignment to balanced randomization and studies uniform performance over finite baseline supports with arbitrary cell masses and unrestricted cellwise arrival probabilities and outcome means, subject to the restrictions in \cref{obj:def:law-class}. This comparison separates identification and regular efficiency from the uniform estimation difficulty generated by many sparsely observed cells.
% [Kallus and Mao](https://academic.oup.com/jrsssb/article/87/2/480/7829031)

A closely related dimension effect appears in \citet[Theorems~1--2 and Appendix~C.5]{ZengBalakrishnanHanKennedy2024}. Their experiment supplies complete covariate--treatment--outcome records and imposes treatment positivity with an unknown categorywise propensity score. Theorem~1 bounds squared error by a sampling term and a dimension term proportional to the squared ratio of alphabet size to sample size. For alphabet sizes at most a constant multiple of sample size times its logarithm, Theorem~2 establishes a lower bound with a squared logarithmic improvement in the dimension term, with constants depending on the fixed positivity level. Appendix~C.5 develops the moment-matching comparison through approximately normalized masses in an auxiliary Poisson experiment and a reduction to probability laws. Here, treatment probabilities are known by design, while outcome arrival varies across observed cells. The arrival-and-dimension characterization in \cref{obj:thm:point-frontier} therefore concerns a different source of incomplete information.
% [Zeng, Balakrishnan, Han, and Kennedy](https://arxiv.org/html/2405.00118v3)

For the complete-record discrete-confounding experiment, \citet[Sharp fixed-interior minimax theorem]{CausalSmith2026} obtains, for each fixed overlap level \(\epsilon\in(0,1/2)\),
\[
\frac1n+\frac{d^2}{n^2(\log n)^2}
\]
up to \(\epsilon\)-dependent constants whenever \(n\ge N_\epsilon\) and \(d\le\rho_\epsilon n\log n\). The same work gives the parametric endpoint rate \(n^{-1}\) under balanced treatment for every positive \(n,d\). These results supply direct precedents for polynomial correction in sparse cells and dimension-independent estimation under complete observation and balanced assignment. In the present experiment, \cref{obj:prop:complete-arrival-reduction} identifies the corresponding complete-arrival reduction, while the main characterization tracks the additional difficulty associated with outcome arrival. In the independent annotation experiment, \citet[Sharp annotation frontier theorem]{CausalSmith2026Annotation} combines \(n\) complete outcome-labeled records with \(m\) independent treatment--covariate records and obtains
\[
\min\left\{1,\frac1n+\frac{d^2}{(n+m)^2\log^2(en)}\right\}
\]
for fixed interior overlap, all \(n\ge1\), \(m\ge0\), and \(d\ge2\), with constants depending on the overlap level. The current experiment instead observes arrival indicators and surrogate measurements within one randomized sample, with cellwise arrival probabilities bounded below by \(q\); its extra term is therefore weighted by \((1-q)^2\) and uses the within-sample arrival information.
% [CausalSmith: fixed-overlap estimation](https://causalsmith.org/papers/stat_discrete_ate_minimax_loggap_polynomial_upper_match/), [CausalSmith: independent annotation](https://causalsmith.org/papers/stat_semisupervised_discrete_ate_annotation_frontier_v1/slides/)

Structured decaying-labeling theory provides a complementary benchmark. \citet[Theorem~2.2 and equation~(2.4)]{ZhangChakraborttyBradic2025} establishes asymptotic normality at an effective sample size determined by the inverse moment of the treatment-and-labeling propensity. Its conditions include correct specification of both nuisance limits, a product-rate restriction on nuisance estimation errors, diverging effective sample size, and the stated moment and tail conditions. Their Theorems~4.1--4.2 develop bias-reducing procedures under offset labeling models, design and noise regularity, and alternative sparsity restrictions. The decoupled procedure additionally controls estimation of the treatment propensity; centered asymptotic normality retains the required control of the product of nuisance approximation errors. These results accommodate declining labeling probabilities through structured nuisance estimation. The present uniform analysis concerns arrival floors between one half and one and freely varying finite-cell nuisance functions. Its dimension parameter measures the number of baseline categories, whereas the sparsity conditions in those structured procedures regulate the complexity of nuisance representations.
% [Zhang, Chakrabortty, and Bradic](https://arxiv.org/html/2305.12789)

Uniform inference with high-dimensional missing data provides an earlier point of contact. \citet[pp.~285--319]{Robins1997} studies how selection probabilities and growing nuisance complexity can obstruct uniform inference in missing-data experiments. Higher-order influence-function analyses subsequently derive minimax conclusions under smoothness, density, and related regularity restrictions \citep[pp.~335--421]{Robins2008}; \citet[pp.~1305--1321]{Robins2009} and \citet[pp.~1951--1987]{Robins2017} develop related minimax results. The present experiment replaces those smooth-function restrictions with a finite but unrestricted collection of cell masses, arrival probabilities, and outcome means. Its contribution is an exact finite-sample arrival-and-dimension rate and a normalized-prior comparison on the complete observed-record space, rather than a higher-order expansion over a smooth nuisance class.

The estimation and converse methods have substantial precedents in discrete functional estimation. \citet[Sections~III--IV]{JiaoVenkatHanWeissman2015} combines polynomial approximation with unbiased estimation of polynomial terms and develops lower bounds through prior comparisons. \citet[Proposition~3 and Appendix~B.2]{WuYang2016} uses moment matching and comparison of Poisson mixtures to establish the dimension-dependent entropy lower bound. These methods explain how aggregate functional estimation can exploit sparse categories even when individual cell parameters are difficult to estimate. Their role here is to guide the polynomial correction for missing outcomes and the construction of statistically close alternatives with separated treatment effects.
% [Jiao, Venkat, Han, and Weissman](https://arxiv.org/pdf/1406.6956), [Wu and Yang](https://arxiv.org/pdf/1407.0381)

The expected-length criterion places inference beside the theory of honest confidence intervals. \citet[Theorem~1 and Section~5]{CaiLow2004} relates uniform coverage and expected length to statistical distinguishability for linear functionals in Gaussian experiments, including nonconvex parameter spaces formed from finite unions of convex classes. In causal inference, \citet{Armstrong2021} studies average treatment effects conditional on realized treatments and covariates under smoothness or shape restrictions: finite-sample optimality uses Gaussian errors with known variance, and feasible procedures with unknown error distributions receive asymptotic guarantees. For randomized trials with missing outcomes, \citet[Theorem~2 under Conditions~2--3]{DiazVanDerLaan2017} establishes doubly robust asymptotic normality under a Donsker condition and nuisance convergence faster than the inverse fourth root of sample size, with at least one correct nuisance limit. The criterion in \cref{obj:thm:interval-frontier} is finite-sample uniform coverage together with worst-law expected length over the unrestricted finite-cell class.
% [Cai and Low](https://arxiv.org/pdf/math/0503662), [Armstrong and Kolesár](https://onlinelibrary.wiley.com/doi/abs/10.3982/ECTA16907), [Díaz and van der Laan](https://arxiv.org/html/1704.01538)

Against these benchmarks, the contribution is the joint characterization of outcome arrival and baseline dimension for minimax squared error and honest expected length in \cref{obj:thm:point-frontier,obj:thm:interval-frontier}. The converse in \cref{obj:thm:normalized-converse} constructs exactly normalized causal laws and compares the induced mixtures of the complete observed record, including baseline category, treatment, surrogate, arrival indicator, and observed outcome mark. This experiment-specific comparison supports both statistical criteria while retaining the information carried by every observed coordinate.
